% !TEX root = ../main.tex
\PlanPage{试点与转化计划}{武桐宇}{38}
\textbf{本页结论：}试点采用 12 周的阶段门设计：先验证资料和规则能否复现，再验证任务价值，最后才讨论付费转化。项目目前没有已签试点协议；以下为执行计划。

阶段门设置的目的不是追求快速成交，而是防止在来源、权限和规则仍不稳定时向外扩展。每一阶段都有可交付物和退出条件；未达成时，优先修正证据/数据问题或停止该场景，而不是用市场叙事覆盖技术缺口。

\begin{table}[htbp]
\centering
\caption{12周试点与转化阶段门（规划）}
\label{tab:wty-pilot-plan}
\small
\begin{tabularx}{\linewidth}{p{0.16\linewidth}p{0.19\linewidth}X p{0.21\linewidth}}
\toprule
阶段 & 周期 & 主要交付与责任 & 继续/退出条件 \\
\midrule
准备 & 第1--2周 & 固定公开案例、版本、权限说明与操作手册 & 来源可打开、样例可复算，否则先修复 \\
任务试用 & 第3--6周 & 研究任务、问题日志、导出记录；团队负责支持 & 能完成来源定位与状态区分，否则调整交互 \\
复核迭代 & 第7--9周 & 第二位复核者复跑，记录差异与修订 & 差异可解释且不涉及未授权数据，否则暂停 \\
转化评估 & 第10--12周 & 交付范围、成本假设与是否付费的书面确认 & 只有书面需求与合规条件齐备才讨论收费 \\
\bottomrule
\end{tabularx}
\end{table}

试点中产生的反馈、日志和个人信息仅在明确同意的范围内使用，汇总展示时去标识化。任何“转化率”“满意度”或“续费率”都必须在实际测量后才可写入计划书。
